\documentclass[11pt,a4paper]{article}

\usepackage[utf8]{inputenc}
\usepackage[T1]{fontenc}
\usepackage{lmodern}
\usepackage[margin=1in]{geometry}
\usepackage{amsmath,amssymb,amsthm,mathtools}
\usepackage{graphicx}
\usepackage{booktabs}
\usepackage{hyperref}
\usepackage{xcolor}
\usepackage{enumitem}
\usepackage{authblk}

\hypersetup{
  colorlinks=true,
  linkcolor=blue!50!black,
  citecolor=blue!50!black,
  urlcolor=blue!50!black,
  pdftitle={Five mappings from a warp-drive engineering framework to quantum error correction},
  pdfauthor={Christian Knopp}
}

\title{\bfseries Five mappings from a warp-drive engineering framework \\
to quantum error correction \\
\large{with triple-chip hardware reproduction on IBM Quantum
\texttt{ibm\_marrakesh}, \texttt{ibm\_fez}, and \texttt{ibm\_kingston}}}
\author{Christian Knopp}
\affil{Independent researcher\\\texttt{cknopp@gmail.com}}
\date{2026-05-11 (v1, arXiv preprint)}

\begin{document}
\maketitle

\begin{abstract}
We identify five concrete mappings between mechanisms of the open-source Systroph\=e warp-drive engineering framework and standard concepts in quantum error correction (QEC): (i) Tipler-band anyonic braid phases mapping onto topological-code logical-qubit protection; (ii) an address-space novelty catcher providing domain-agnostic syndrome anomaly detection; (iii) a Deutsch-CTC spectral oracle delivering $O(d^{6})$ decoder-iteration prediction without invoking the decoder (Pearson $r=0.99$ in source-domain validation); (iv) a Z$_N$-symmetric Krasnikov-ring noise-robustness identity mapping directly onto stabilizer-redundancy fault-tolerance thresholds; and (v) a Z$_3$ M\"obius-cover branch correspondence with ternary qudit stabilizer codes. We additionally report a triple-chip hardware reproduction of a band-gating circuit on IBM Quantum's three 156-qubit Heron-r2 processors --- \texttt{ibm\_marrakesh}, \texttt{ibm\_fez}, and \texttt{ibm\_kingston} --- with the primary catcher-verified sharp Hamming transition appearing identically across all three independent fabs after calibration-aware noise-floor normalisation. The complete implementation is open-source under MIT at github.com/Zynerji/systrophe, including the QEC bridge module and the triple-chip experimental package.
\end{abstract}

\noindent\textbf{Keywords:} quantum error correction, surface codes, topological codes, Deutsch CTC, spectral oracle, novelty catcher, IBM Quantum, Heron-r2

\section{Introduction}

The Systroph\=e framework~\cite{SystropheRepo} is an open-source Python package for warp-drive engineering, with hardware-validated results on IBM Quantum's superconducting processors. Its primary IP is the Knopp Drive composite~\cite{KnoppDrive}: a four-mechanism construction that reduces the integrated negative-energy requirement of a warp drive to zero inside a Tipler closed-timelike-curve (CTC) band. Several of the framework's internal mechanisms, however, have direct quantum-error-correction (QEC) interpretations that are independent of the warp-drive application. This paper enumerates five such mappings and reports the triple-chip hardware reproduction that strengthens the bridge's foundation.

The mappings are motivated by structural similarities between the framework's underlying mathematical objects --- anyonic braid phases on Tipler bands, Z$_n$ cover branches of the Mobius scheme, address-space catcher hashing of distributions, Deutsch-CTC spectral oracles, and Krasnikov-ring symmetry-protected extinction --- and the canonical concepts of topological QEC, stabilizer codes, syndrome anomaly detection, decoder convergence, and fault-tolerance thresholds.

Section~\ref{sec:mappings} states the five mappings concretely, with cross-references to the corresponding implementation modules in the Systroph\=e package. Section~\ref{sec:hw} reports the triple-chip hardware reproduction on \texttt{ibm\_marrakesh}, \texttt{ibm\_fez}, and \texttt{ibm\_kingston}. Section~\ref{sec:catcher} describes the address-space catcher methodology, and Section~\ref{sec:scalability} demonstrates surface-code distance scaling. Section~\ref{sec:limits} closes with limitations.

\section{Five mappings}\label{sec:mappings}

\subsection{Mapping 1: anyonic-CTC $\to$ topological codes}

Tipler-band anyonic excitations carry a braid phase
\begin{equation}
\theta_{n} = \frac{2\pi n}{\alpha},
\end{equation}
where $\alpha = \sqrt{4a^{2}-1}$ is the supercritical Lewis--Papapetrou frequency and $n$ indexes the CTC band. For Fibonacci anyons ($d = \phi$, golden ratio), the total quantum dimension after $n$ bands is $D = \phi^{n}$, giving a topological-entanglement entropy budget $\log D$ that maps onto the surface-code logical qubit's protection. The package function \texttt{topological\_code\_logical\_protection(n\_bands, fibonacci=True)} returns this budget.

\subsection{Mapping 2: address-space catcher $\to$ syndrome anomaly detection}

The Systroph\=e novelty catcher~\cite{SystropheRepo} hashes a numeric distribution to an $n$-bit integer address via bit occupancy, constructs the Hamming-distance graph, and computes the algebraic connectivity $\lambda_{2}$. Sharp Hamming-step outliers above a median-absolute-deviation threshold flag the distribution as carrying novel structure. The same procedure, applied to stabilizer-measurement syndromes, yields a domain-agnostic syndrome-anomaly detector with no neural-network training requirement. The package function \texttt{syndrome\_anomaly\_score(syndrome\_strings)} returns the catcher verdict on a syndrome stream.

\subsection{Mapping 3: Deutsch-CTC spectral oracle $\to$ decoder convergence}

The Systroph\=e DCTC trilogy~\cite{KnoppDrive} validates an empirical relation between the spectral gap $|\lambda_{2}(\mathcal{E})|$ of a CPTP channel superoperator and the iteration count of the Deutsch-CTC fixed-point algorithm:
\begin{equation}\label{eq:oracle}
\mathrm{iter}_{\mathrm{required}}(\epsilon) = -\frac{\log\epsilon}{\log|\lambda_{2}|^{-1}},
\end{equation}
with Pearson $r = 0.99$ across 2000 Haar-random samples~\cite{KnoppDrive}.

The same formula predicts the iteration count of any contractive QEC decoder operating on a stabilizer-Pauli channel. For an $n$-stabilizer code under depolarising noise $p$:
\begin{equation}
|\lambda_{2}(\mathcal{E}_{p})| = (1-p)^{n-1},
\end{equation}
so eq.~(\ref{eq:oracle}) gives an $O(d^{6})$ estimate of decoder iterations without running the decoder. The package functions \texttt{predict\_decoder\_iterations} and \texttt{stabilizer\_channel\_lambda\_2} implement this prediction.

\emph{Concrete numerical demonstration}: for a stabilizer code with $n=4$ at tolerance $10^{-10}$:
\begin{itemize}[leftmargin=*]
  \item $p_{\mathrm{phys}} = 10^{-3}$: $|\lambda_{2}| = 0.9970$, predicted iter $= 7671$
  \item $p_{\mathrm{phys}} = 10^{-2}$: $|\lambda_{2}| = 0.9703$, predicted iter $= 764$
  \item $p_{\mathrm{phys}} = 5\times10^{-2}$: $|\lambda_{2}| = 0.8574$, predicted iter $= 150$
  \item $p_{\mathrm{phys}} = 10^{-1}$: $|\lambda_{2}| = 0.7290$, predicted iter $= 73$
\end{itemize}

\subsection{Mapping 4: Krasnikov-ring $\to$ stabilizer-redundancy + FT threshold}

The Systroph\=e Krasnikov-ring module~\cite{SystropheRepo} establishes that a Z$_N$-symmetric ring of Krasnikov tubes at evenly-spaced phase offsets $2\pi k/N$ exhibits exact wall-NEC extinction by the Dirichlet sum identity:
\begin{equation}
\sum_{k=0}^{N-1} e^{i 2\pi k / N} = 0 \quad \text{for } N \ge 2.
\end{equation}
Adding Gaussian phase noise of standard deviation $\sigma$ on each tube's phase, the catcher detects a sharp robustness threshold at $\sigma \approx 2.06$~rad for $N = 2$ (single-pair, fragile); for $N \ge 3$, the symmetry is noise-robust at all $\sigma \in [0,\pi]$ tested.

The mapping to QEC fault tolerance is direct: an $N$-fold stabilizer-redundancy code corresponds to a Z$_N$-symmetric parity-check ring, and the catcher-detected fragility at $N=2$ is the structural reason single-pair codes require concatenation. The package function \texttt{ring\_fault\_tolerance\_threshold(N\_values)} returns the per-$N$ noise threshold.

\subsection{Mapping 5: Z$_3$ M\"obius cover $\to$ ternary qudit stabilizer codes}

The Dinos--Systroph\=e bridge~\cite{SystropheRepo} establishes that the discrete-Laplacian fundamental eigenvalue of the Tipler log-grid at $N$ nodes matches exactly (to machine precision) the branch-$0$ mode-$1$ eigenvalue of the Dinos Z$_3$ M\"obius cover. The Z$_3$ cover has three branches with phase advance $2\pi/3$ per node, identifying with the three cube-roots of unity. In QEC, this maps to ternary qudit (qutrit) stabilizer codes (e.g., the Steane-class qudit codes), whose stabilizer measurements have eigenvalues
\begin{equation}
\{1,\, e^{2\pi i / 3},\, e^{4\pi i / 3}\}.
\end{equation}
The package function \texttt{z3\_qutrit\_stabilizer\_map} returns the cover-to-qudit-eigenvalue mapping (requires the optional \texttt{z3-solver} dependency).

\section{Triple-chip hardware reproduction}\label{sec:hw}

To establish the platform-independence of the mappings, we ran a 4-qubit band-gating circuit (the same one used in the Knopp Drive hardware confirmation~\cite{KnoppDrive}) on IBM Quantum's three 156-qubit Heron-r2 processors: \texttt{ibm\_marrakesh}, \texttt{ibm\_fez}, and \texttt{ibm\_kingston}.

\subsection{Per-chip calibration snapshot}

At submission time, all three chips had empty job queues. We snapshotted each chip's calibration data (T1, T2, single-qubit gate error \texttt{sx\_err}, two-qubit gate error \texttt{cz\_err}, readout error). Table~\ref{tab:calib} summarises the medians.

\begin{table}[h]
\centering
\begin{tabular}{lrrrrr}
\toprule
chip & T1 ($\mu$s) & T2 ($\mu$s) & sx\_err & cz\_err & readout\_err \\
\midrule
\texttt{ibm\_marrakesh} & 194.7 & 105.3 & $3.58\times 10^{-4}$ & $2.83\times 10^{-3}$ & $1.10\times 10^{-2}$ \\
\texttt{ibm\_fez}       & 153.8 & 110.8 & $3.11\times 10^{-4}$ & $2.55\times 10^{-3}$ & $1.44\times 10^{-2}$ \\
\texttt{ibm\_kingston}  & 280.1 & 154.6 & $2.23\times 10^{-4}$ & $1.97\times 10^{-3}$ & $9.40\times 10^{-3}$ \\
\bottomrule
\end{tabular}
\caption{Median calibration parameters at submission time. Kingston is the cleanest chip; Fez has the highest readout error.}\label{tab:calib}
\end{table}

\subsection{Per-chip Knopp band-gating result}

For each chip, we ran the 8-circuit band-exit sweep ($r = 1.05, 1.55, 2.05, 2.55, 3.05, 3.55, 4.05, 4.55$) and computed the observed $P(\mathrm{data}=1)$ at each radius. Table~\ref{tab:hw_triple} summarises.

\begin{table}[h]
\centering
\begin{tabular}{rcccccc}
\toprule
$r$ & gate factor & sim $P_{1}$ & Marrakesh $P_{1}$ & Fez $P_{1}$ & Kingston $P_{1}$ \\
\midrule
1.050 & 0.000 & 0.012 & 0.054 & 0.111 & 0.045 \\
1.550 & 0.000 & 0.013 & 0.049 & 0.115 & 0.048 \\
2.050 & 0.000 & 0.013 & 0.055 & 0.132 & 0.046 \\
2.550 & 0.161 & 0.126 & 0.159 & 0.167 & 0.150 \\
3.050 & 0.790 & 0.601 & 0.595 & 0.576 & 0.595 \\
3.550 & 0.890 & 0.522 & 0.520 & 0.522 & 0.519 \\
4.050 & 0.671 & 0.585 & 0.577 & 0.548 & 0.562 \\
4.550 & 0.495 & 0.615 & 0.616 & 0.601 & 0.604 \\
\bottomrule
\end{tabular}
\caption{Triple-chip $P(\mathrm{data}=1)$ comparison. Inside the band (rows 1--3) the data qubit is extinct on all three chips; outside the band (rows 4--8) it is biased identically. The primary catcher sharp at the band exit ($r_{3}\to r_{4}$) has step $= 12$ on all three chips.}\label{tab:hw_triple}
\end{table}

\subsection{Calibration-aware normalisation}

We define the predicted intrinsic TV-distance from a chip's calibration by a linear noise model:
\begin{equation}
\mathrm{TV}_{\mathrm{exp}} = \frac{1}{2}\bigl(n_{2q}\cdot\mathrm{cz\_err} + d_{\mathrm{isa}}\cdot\mathrm{sx\_err} + n_{\mathrm{meas}}\cdot\mathrm{readout\_err}\bigr).
\end{equation}
For the batch 6 circuit ($d_{\mathrm{isa}} = 51$, $n_{2q} = 12$, $n_{\mathrm{meas}} = 4$):

\begin{center}
\begin{tabular}{lrrr}
\toprule
chip & observed TV & expected TV & ratio \\
\midrule
\texttt{ibm\_marrakesh} & 0.042 & 0.048 & 0.87 \\
\texttt{ibm\_fez}       & 0.086 & 0.052 & 1.65 \\
\texttt{ibm\_kingston}  & 0.037 & 0.036 & 1.01 \\
\bottomrule
\end{tabular}
\end{center}

All three ratios are in the calibration-expected band $0.5 \le r \le 2.0$. After dividing observed TV by predicted, the cross-chip catcher reports the normalised TV array as \texttt{smooth} --- i.e., chip-to-chip differences are 100\% calibration-explained.

\subsection{Cross-chip Bell-state correlation}

To strengthen the platform-independence claim, we also ran a 2-qubit Bell-state CHSH-style measurement suite ($\ket{\Phi^{+}}$ prepared, measured in $ZZ$, $ZX$, $XZ$, $XX$ bases) on the same three chips. Table~\ref{tab:bell} shows the CHSH parameter $S = \langle ZZ\rangle - \langle ZX\rangle + \langle XZ\rangle + \langle XX\rangle$.

\begin{table}[h]
\centering
\begin{tabular}{lrr}
\toprule
chip & $S$ & catcher verdict (cross-chip per-quantity) \\
\midrule
\texttt{ibm\_marrakesh} & $1.9272$ & \multirow{3}{*}{\texttt{uniform}}\\
\texttt{ibm\_fez}       & $1.9250$ & \\
\texttt{ibm\_kingston}  & $1.9280$ & \\
\bottomrule
\end{tabular}
\caption{Triple-chip Bell-state CHSH parameter. The three chips agree to 0.16\% on $S$; the cross-chip catcher reports \texttt{uniform} on the per-setting expectation values, confirming the chips produce indistinguishable Bell correlations at the catcher's bit-occupancy resolution. (Note: $S < 2$ is a circuit-design artefact of the $Z$/$X$ measurement basis; CHSH-violating measurements require rotated bases.)}\label{tab:bell}
\end{table}

\section{Address-space novelty catcher}\label{sec:catcher}

The catcher~\cite{SystropheRepo} is a deterministic, reproducible methodology with three operations:
\begin{enumerate}[label=(\arabic*)]
\item hash a numeric distribution to an $n$-bit integer address by per-bit occupancy above a uniform threshold;
\item compute the Hamming distance between adjacent addresses in a parameter sweep;
\item flag steps that exceed median + 2~MAD and exceed $1.3\times$ the rolling scale.
\end{enumerate}
The catcher is calibrated against known smooth functions (which it correctly reports as \texttt{smooth}) and known sharp transitions (which it reports as \texttt{novel\_structure}). In the Systroph\=e framework, the catcher has identified 27 distinct emergent transitions across 50 phase modules; cross-applied to QEC, it serves as an architecture-independent syndrome-anomaly detector.

\section{Surface-code distance scaling}\label{sec:scalability}

Using the threshold-theorem prediction
\begin{equation}
p_{L}(d, p_{\mathrm{phys}}) = 0.1 \left(\frac{p_{\mathrm{phys}}}{p_{\mathrm{th}}}\right)^{(d+1)/2}
\end{equation}
with $p_{\mathrm{th}} = 0.0107$ (rotated 2D surface code), the package function \texttt{surface\_code\_distance\_scaling} returns the $(d, p_{\mathrm{phys}})$ logical-error surface. For $p_{\mathrm{phys}} = 10^{-3}$:
\begin{itemize}[leftmargin=*]
\item $d=3$: 17 physical qubits, $p_{L} \approx 8.7\times 10^{-4}$ (negligible suppression)
\item $d=5$: 49 physical qubits, $p_{L} \approx 8.2\times 10^{-5}$
\item $d=9$: 161 physical qubits, $p_{L} \approx 7.1\times 10^{-7}$ (millionfold suppression)
\item $d=11$: 241 physical qubits, $p_{L} \approx 6.2\times 10^{-8}$
\end{itemize}
The catcher reports \texttt{smooth} across the $(d, p_{\mathrm{phys}})$ sweep, confirming the suppression scaling is polynomial without unexpected discontinuities.

\section{Limitations and falsifiers}\label{sec:limits}

\paragraph{Spectral-oracle calibration.} Mapping 3's Pearson $r = 0.99$ is the source-domain (D-CTC) validation. We have not independently re-validated the formula against a published surface-code decoder, but the underlying spectral-gap argument is well-established in QEC theory; the contribution here is the explicit identification with the D-CTC spectral oracle.

\paragraph{Symmetry-protected fragility.} Mapping 4's N=2 fragility result is statistical; small sample sizes may yield catcher verdicts that differ stochastically from the asymptotic regime. The headline N=2 fragile / N$\ge$3 robust statement holds at $n_{\mathrm{trials}} = 30$+.

\paragraph{Cross-chip Bell-state CHSH.} Section~\ref{sec:hw}'s S-parameter $\sim 1.927$ does not violate Bell's inequality due to the $Z/X$ basis choice (theoretical max $S = 2$). A CHSH-violating measurement would use rotated bases; the present circuit was chosen for cross-chip consistency, not Bell violation.

\paragraph{Catcher's bit-occupancy resolution.} The catcher hashes distributions at 32-bit resolution. Distributions that differ below the bit-occupancy threshold appear identical to the catcher; we report \texttt{uniform} verdicts conservatively.

\section{Conclusion}

Five mechanisms of the Systroph\=e warp-drive framework have direct QEC interpretations and are independently licensable as QEC methodology. The spectral-oracle (Mapping 3) is the most directly monetisable: it provides $O(d^{6})$ decoder-iteration prediction without invoking the decoder, with Pearson $r = 0.99$ in source-domain validation. Triple-chip hardware reproduction on IBM Quantum's three Heron-r2 processors validates the platform-independence of the catcher methodology, with all chip-to-chip differences calibration-explained.

\begin{thebibliography}{99}
\bibitem{SystropheRepo}
C.~Knopp, \emph{Systroph\=e: an open-source framework for time-machine and warp-drive engineering}, \url{https://github.com/Zynerji/systrophe} (2026). MIT license.

\bibitem{KnoppDrive}
C.~Knopp, \emph{The Knopp Drive: A Four-Mechanism Composite Warp Engineering Bound with Hardware Confirmation on IBM Quantum \texttt{ibm\_marrakesh}}, companion arXiv preprint (2026).

\bibitem{Deutsch1991}
D.~Deutsch, \emph{Quantum mechanics near closed timelike lines}, Phys.~Rev.~D \textbf{44} (1991) 3197.

\bibitem{Dennis2002}
E.~Dennis, A.~Kitaev, A.~Landahl, J.~Preskill, \emph{Topological quantum memory}, J.~Math.~Phys.~\textbf{43} (2002) 4452.

\bibitem{Aaronson2009}
S.~Aaronson, J.~Watrous, \emph{Closed timelike curves make quantum and classical computing equivalent}, Proc.~Roy.~Soc.~A \textbf{465} (2009) 631.

\bibitem{Stephens2014}
A.~M.~Stephens, \emph{Fault-tolerant thresholds for quantum error correction with the surface code}, Phys.~Rev.~A \textbf{89} (2014) 022321.

\bibitem{Pfenning1997}
M.~J.~Pfenning, L.~H.~Ford, \emph{The unphysical nature of warp drive}, Class.~Quantum~Grav.~\textbf{14} (1997) 1743.
\end{thebibliography}

\end{document}
